\chapter*{Course Map: From Packets to WANs}
\addcontentsline{toc}{chapter}{Course Map: From Packets to WANs}\markboth{Course Map}{}
The labs progress from packet observation to device configuration, routing, services, Layer-2 switching, and WAN protocols.

\begin{center}
\begin{tikzpicture}[node distance=8mm,stage/.style={draw=MLBlue,fill=MLPaleBlue,rounded corners=2mm,text width=13cm,minimum height=1.45cm,align=left,inner sep=3mm,font=\small\color{MLNavy}},arr/.style={-{Stealth[length=2mm]},thick,draw=MLTeal}]
\node[stage] (a) {\textbf{Stage 1 -- Observe and configure: Labs 1--2}\\Wireshark captures, eNSP topology building, VRP navigation, interface addressing, verification, and saving.};
\node[stage,below=of a] (b) {\textbf{Stage 2 -- Make networks reachable: Labs 3--4}\\Static and backup routes, single-area OSPF, routing tables, neighbors, and connectivity tests.};
\node[stage,below=of b] (c) {\textbf{Stage 3 -- Add network services: Labs 5--6}\\FTP with AAA authorization; DHCP global and interface pools, leases, and DNS options.};
\node[stage,below=of c] (d) {\textbf{Stage 4 -- Build the switched LAN: Labs 7--8}\\Manual aggregation, LACP priorities, VLAN port types, PVIDs, and isolation tests.};
\node[stage,below=of d] (e) {\textbf{Stage 5 -- Connect sites across a WAN: Lab 9}\\HDLC, PPP, DCE clocking, OSPF, PAP/CHAP authentication, and debugging.};
\node[stage,below=of e,draw=MLGreen,fill=MLPaleTeal] (f) {\textbf{Group project -- Integrate and explain}\\Build and demonstrate an eNSP network, justify protocol choices, explain forwarding, and report the evidence.};
\draw[arr] (a)--(b);\draw[arr] (b)--(c);\draw[arr] (c)--(d);\draw[arr] (d)--(e);\draw[arr] (e)--(f);
\end{tikzpicture}
\end{center}

\section*{How to use these notes}
Each chapter presents the source material in four forms:
\begin{itemize}
\item \textbf{Source visuals} show key topologies, interfaces, and protocols.
\item \textbf{Command blocks} preserve the Huawei VRP configuration sequence.
\item \textbf{Verification evidence} identifies the state, capture, or test that proves success.
\item \textbf{Troubleshooting order} provides a repeatable diagnostic workflow.
\end{itemize}

\begin{keyidea}
Use the same loop in every lab: \textbf{build $\rightarrow$ configure $\rightarrow$ inspect $\rightarrow$ test $\rightarrow$ explain}. Only the protocol changes.
\end{keyidea}
